\chapter{What scaling laws don't tell you}
\label{ch:what-scaling-omits}

\newthought{Both Kaplan and Chinchilla} are about \emph{training} cost. Their
loss curves assume that the model, once trained, can be \emph{used} at the
context length it was trained at. That assumption was nearly free in 2020.
It is the bottleneck in 2026.

\section{Three things the laws are silent on}

\paragraph{1. Inference cost as a function of context length.} The Kaplan
power law says that loss falls smoothly in parameters \(N\), data \(D\), and
training compute \(C\). It does not say a word about how much \emph{compute
per query} a trained model needs at sequence length \(n\). That compute
grows as \(\bigO(n^2)\) for dense attention (Chapter~\ref{ch:why-quadratic}),
and that growth has no analog on the scaling-laws plot.

\paragraph{2. The KV cache.} The cache size grows as \(\bigO(n)\) per layer.
For long contexts and large models, it can exceed the model's weight footprint
by a wide margin. Scaling laws are written about a model that fits on one
GPU and serves one short request at a time; they have nothing to say about
a model whose KV cache for a single request requires multi-GPU sharding.

\paragraph{3. The quality--length curve.} A nominal context window says how
many tokens a model accepts. A \emph{functional} context window says how
many tokens a model can actually reason over. There is no scaling law that
relates the two. The community-standard benchmarks
(\citealp{hsieh2024ruler, chen2025mrcr}) consistently find that nominal
windows of 128K tokens have functional windows much smaller, and that the
gap depends on architecture in ways that scaling laws do not capture.

\section{Why the agentic workloads broke the model}

Through about 2022 the typical LLM workload was: short user query
(${\sim}100$ tokens), short response (${\sim}500$ tokens). Total context
${\sim}600$ tokens. For that regime, dense attention's \(n^2\) is small,
the KV cache fits on a single card, scaling laws describe the behaviour you
care about, and the world is fine.

The 2023--2026 wave broke that regime. The new workloads:

\begin{itemize}
\item A coding agent reading an entire 400k-token monorepo and refactoring
across files.
\item A legal assistant reading a 200-page contract and resolving
cross-references.
\item A research summariser reading 50 PDFs and producing a structured
literature map.
\item A long-running agent session whose ``context'' is the trace of every
prior tool call, plan, and observation.
\end{itemize}

These do not look like 600-token chats. They look like 100k--1M-token
contexts, served at frontier-model quality. And the scaling-laws-era
infrastructure — dense transformer, single H100, FlashAttention,
GQA — does not deliver them at any reasonable price.

\section{Two responses, and the rest of the book}

Faced with the gap between what scaling laws describe (training compute)
and what 2026's products need (cheap, accurate long-context
inference), the field has split into two camps.

\begin{description}
\item[\textnormal{\textbf{Camp 1: don't change the algorithm, change the
implementation.}}] FlashAttention (Chapter~\ref{ch:flashattention}),
better KV-cache schemes (Chapter~\ref{ch:kv-cache}), MQA / GQA, paged
caches, sliding windows. The asymptote stays \(n^2\); the constants and
the memory bills get much better. This is where most production
deployments live in 2026.

\item[\textnormal{\textbf{Camp 2: change the algorithm.}}] State-space models
(Chapter~\ref{ch:ssm}), long convolutions and linear attention
(Chapter~\ref{ch:linattn}), hybrids (Chapter~\ref{ch:hybrids}), and
content-dependent sparse mechanisms like SSA
(Chapter~\ref{ch:ssa}). The asymptote becomes
\(\bigO(n)\) or \(\bigO(n \log n)\); the trade-off is that something else
about the model has to be different — finite state, fixed sparsity
patterns, or proprietary selection routing.
\end{description}

\keyidea{Scaling laws describe training-loss returns on training resources.
They do not describe inference-time cost as a function of context length,
and that gap is the entire reason the post-transformer landscape exists.}

\section{Where the rest of the book goes}

Part~IV (next) is Camp 1. Part~V is Camp 2. Part~VI is the 2$\times$2
frame that makes the whole landscape make sense, and the careful read of
SubQ's bet inside it.

\fillme{FILL-09-01}{Sidebar: ``lost in the middle''}{%
  ${\sim}300$-word sidebar on Liu et al.\ 2023 ``Lost in the Middle: How
  Language Models Use Long Contexts''. Convey: nominal context window is
  not the same as functional context; LLMs put more attention on the
  beginning and end of the prompt; prompt-engineering implication.
  Connect to the RULER and MRCR benchmarks introduced later.}{%
  Liu et al.\ 2023; Hsieh et al.\ 2024 (RULER).}{%
  ${\sim}300$ words; no equations.}

\fillme{FILL-09-02}{Worked example: 1M-token agent TCO over 1B requests}{%
  Total cost of ownership for a hypothetical agent that serves 1B sessions
  per year, each averaging 200k tokens of context plus 5k tokens of
  generation. Assume Llama-3 70B class model on H100. Show: serving FLOPs,
  KV-cache memory pressure, GPU-hours, and \$ per session at on-demand
  pricing. Then re-do the row for a 10$\times$ cheaper architecture
  (foreshadow Part~V). End with one sentence on what this implies for
  pricing pages of frontier APIs in 2026.}{%
  KV-cache equation from Chapter~\ref{ch:kv-cache}; current cloud pricing.}{%
  Half a page; 2-row table + 3 sentences.}
